\label{sec:rerun}

The census frame is fixed at the snapshot of 19 July 2026, so re-running it probes
the same servers, at the same declared versions, with the same launch commands. What
those commands resolve is not fixed. Each launch installs the package and whatever
its dependency ranges admit on the day. We re-ran the complete frame on \SeptDate{},
sixty-nine days later, with the harness unchanged.

\paragraph{What the re-run found}
\NSeptHandshake{} of \NSeptGraded{} graded servers complete a handshake,
\SeptHandshakeRate{} against \HandshakeRatePoint{} in July. \NLost{} servers that
answered in July no longer answer, and \NGained{} that did not answer now do.

Conformance did not change. On the \NPaired{} servers that answered in both runs,
error-as-result was \PairedEarBefore{} in July and \PairedEarAfter{} in September, silent
type non-enforcement \PairedTypeBefore{} and \PairedTypeAfter{}, and failure to recover
from a malformed frame \PairedMalBefore{} and \PairedMalAfter{}. Rates computed over each
run's own responders move slightly more, type non-enforcement from
\NoTypecheckRateCl{} to \SeptNoTypecheck{}, but that is a change in who answers rather
than in how they answer: the servers September lost had a type non-enforcement rate of
\LostTypeRate{} in July, so removing them raises the rate among those that remain. The
instrument reports the same behavior from the same servers. What changed is how many of
them start.

\paragraph{One dependency release accounts for most of the loss}
Of the \NLost{} servers that stopped answering, \NLostUnchanged{} were probed at the
same declared version with the same launch command as in July, so the servers
themselves did not change. \NSDKTwoBroke{} of them fail with one of two signatures in
their own output: the official Python SDK's compatibility message, which states that
this is \texttt{mcp} 2.x and that \texttt{FastMCP} has been renamed to
\texttt{MCPServer}, or an \texttt{AttributeError} for a handler-registration
decorator that the same major version removed from the low-level
\texttt{Server} class.

That SDK published 2.0.0 on \SDKTwoDate{}, nine days after our snapshot.
\NSDKTwoRange{} of the \NSDKTwoBroke{} declare the dependency as a floor or a range
rather than an exact pin, so the same launch command resolved 1.x in July and 2.x in
September.

Three observations separate this from a fault in our own instrument or host.
The effect is confined to one package ecosystem: \NpmLost{} of \NpmWas{} npm servers
that answered in July stopped answering (\NpmLostRate{}), against \PypiLost{} of
\PypiWas{} on PyPI (\PypiLostRate{}). A host, network or harness fault would not
respect that boundary. The attribution comes from the servers' own stderr rather than
from our classification. And the failure is reversible: we re-launched a random
\NPinTested{} of the affected servers with the SDK constrained below 2.0 and
otherwise unaltered, and \NPinRestored{} completed a handshake.

Restricting attention to the servers this could affect, \NPySDKWas{} PyPI servers
attributed to a Python MCP SDK completed a handshake in July and \NPySDKNow{} do now.
\PySDKLostCI{} of them stopped working, and no author changed anything.

\paragraph{What this run is and is not}
It is not an independent census. It is the same frame re-probed, so it says nothing
about servers published after July, and the \NGained{} recoveries are not a sample of
anything. The July run remains the measurement this paper reports, because its
conformance verdicts were collected when the population was largest. The re-run's
contribution is the comparison.
